%%%%%%%%%%%%%%%%%%%%%%%%%%%%%%%%%%%%%%%%%%%%%%%%%%%%%%%%%%%%%%%%%%
%                           Document style
%%%%%%%%%%%%%%%%%%%%%%%%%%%%%%%%%%%%%%%%%%%%%%%%%%%%%%%%%%%%%%%%%%
\documentclass[12pt]{article}

% Importing report style and the first pages
%%%%%%%%%%%%%%%%%%%%%%%%%%%%%%%%%%%%%%%%%%%%%%%%%%%%%%%%%%%%%%%%%%
%                       Main document style
%%%%%%%%%%%%%%%%%%%%%%%%%%%%%%%%%%%%%%%%%%%%%%%%%%%%%%%%%%%%%%%%%%

% Importing user packages (librarys for different functions and commands)
\usepackage[utf8]{inputenc}
\usepackage{subfiles}               % add subfiles
\usepackage[T1]{fontenc}            % text font
\usepackage{lmodern}                % font
\usepackage{pdfpages, pdflscape}    % add pdf's
\usepackage{gensymb, textcomp}      % add symbols
\usepackage{lastpage}               % get the last page number
\usepackage{graphicx}               % add images
\usepackage{setspace}               % define size of space (infinite)
\usepackage{xcolor}                 % colors
\usepackage{comment}                % comments
\usepackage{fancyhdr}               % to get fancy headers
\usepackage [background=black, arrow=red, text=white]{callouts}
\usepackage{wrapfig}                % positioning figures
\usepackage{subcaption}             % position mulitple figures

% page settings
\setlength{\marginparwidth}{2cm}    % margins
\usepackage[a4paper,left=2.5cm,right=2.5cm,top=2.8cm,bottom=4cm]{geometry}
\setlength\parindent{0pt}
\setlength{\parskip}{1em}

\setlength{\headheight}{37pt} % header hight

\cfoot\thepage % page number as footer

% creating base header
\pagestyle{fancy}
\rhead{\includegraphics[height=1.2cm]{images/kir-dev-logo.png}}                   % itt a magasságán lehet finomhangolni

%define colors for links
\definecolor{ceruleanblue}{rgb}{0.16, 0.32, 0.75}
\definecolor{coolblack}{rgb}{0.0, 0.18, 0.39}


% clickable table of contents
\usepackage{hyperref}
\hypersetup{linktoc = all, colorlinks=true,
    linkcolor=coolblack, filecolor=magenta,
    urlcolor=ceruleanblue, citecolor=ceruleanblue}


%%%%%%%%%%%%%%%%%%%%%%%%%%%%%%%%%%%%%%%%%%%%%%%%%%%%%%%%%%%%%%%%%%
%                   Command for the first page
%%%%%%%%%%%%%%%%%%%%%%%%%%%%%%%%%%%%%%%%%%%%%%%%%%%%%%%%%%%%%%%%%%

\newcommand*\FirstPage{
\thispagestyle{empty}
% placing a pdf as background and adding text on top
\includepdf[pages=1,picturecommand*={
    % selecting placement, font size and text type
    \put(160,630){
        \fontsize{28}{28}\selectfont{
            \textsc{\textbf{SZERVEZETI ÉS}}
        }
    }
    \put(85,590){
        \fontsize{28}{28}\selectfont{
            \textsc{\textbf{MŰKÖDÉSI SZABÁLYZAT}}
        }
    }
    \put(250,540){                                              % Itt az X koordinátát lehet finomhangolni
        \fontsize{28}{28}\selectfont{
            \Large{\textbf{Kir-Dev}}                            % Itt nevezd át a körödre
        }
    }
    \put(55,70){
        \fontsize{28}{28}\selectfont{
            \Large{\textbf{Budapest, 2024. október 1.}}
        }
    }
}]{images/kir-dev-frontpage-bg.pdf}}


% section header
\lhead{Szervezeti és Működési Szabályzat - Kir-Dev}
\linespread{1.25}

%%%%%%%%%%%%%%%%%%%%%%%%%%%%%%%%%%%%%%%%%%%%%%%%%%%%%%%%%%%%%%%%%%
%                       New Commands
%%%%%%%%%%%%%%%%%%%%%%%%%%%%%%%%%%%%%%%%%%%%%%%%%%%%%%%%%%%%%%%%%%

%creating paragraphs --> \newpara
\newcommand{\newpara}{\vskip 2mm}

%add pdf's --> \addpdf
\newcommand*\addpdf[1]{
\includepdf[pages=-]{#1}}

%add pdf's in landscape --> \addpdflandscape
\newcommand*\addpdflscape[1]{
\begin{landscape}
\newgeometry{margin=0cm}
\includepdf[pages=-,angle=90]{#1}
\end{landscape}}

%add sections with 4 numbers --> \paragraph
\usepackage{titlesec}
\setcounter{secnumdepth}{4}
\titleformat{\paragraph}
{\normalfont\normalsize\bfseries}{\theparagraph}{1em}{}
\titlespacing*{\paragraph}
{0pt}{3.25ex plus 1ex minus .2ex}{1.5ex plus .2ex}

\titleformat{\section}
  {\normalfont\Large\bfseries}{\thesection.\S}{1em}{}


% sections without numbers, but still in table of contents
%\setcounter{secnumdepth}{0} % remove "%" to include

%=================================================================
%                           Start Document
%=================================================================
\begin{document}

% adding the first two pages
\FirstPage
\pagebreak

A Schönherzes Villamosmérnökök és Informatikusok Egyesületének (továbbiakban: Egyesület) a
Körök alapítására vonatkozó szabályozása alapján a Sysadmin (továbbiakban: Kör)
\linebreak
Szervezeti és
Működési Szabályzatát (továbbiakban: SZMSZ) az alábbiakban állapítja meg:

\section{A Kör alap adatai}

\begin{enumerate}
  \item A Kör neve: Sysadmin
  \item A Kör Reszortja: Kollégiumi Számítástechnikai Kör
  \item A Kör hivatalos elérhetősége: sysadmin@sch.bme.hu
  \item A Kör alapításának éve: 2014
\end{enumerate}

% Section border ==========================================================

\section{A Kör alapvető tevékenysége}

A Kör célja a Schönherz Kollégium (továbbiakban: Ház) és a BME-VIK hallgatói informatikai
szolgáltatásaihoz szükséges számítógépes infrastruktúra karbantartása, fejlesztése és üzemeltetése,
a Kollégiumi Számítástechnikai Kör (továbbiakban: KSZK) többi körével együttműködve elsősorban a
Ház többi egységében felmerülő informatikai igények kielégítése, informatikai tanácsadás nyújtása.

A Kör célja továbbá tagjai szakmai fejlődésének elősegítése, külsős és belsős tanfolyamok,
workshopok tartásával.

% Section border ==========================================================

\section{A Kör felépítése}

\begin{enumerate}
  \item A Kör tagságára jelentkezhet minden rendes Egyesületi tag, aki elfogadja és magára nézve kötelezőnek ismeri el a Kör SZMSZ-ét, a KSZK reszort SZMSZ-ét, valamint egyetért a Kör céljaival és feladataival.
  
  \item A Kör tagja az lehet, aki teljesíti a Reszort felvételi kritériumait.

  \item A Kör posztjai

\begin{enumerate} % posztok
  \item Körvezető
    \begin{enumerate}
    \item Felelősségeiről és feladatköréről a 4. § rendelkezik.
    \end{enumerate}

  \item Vezető Rendszergazda (továbbiakban: VRG)
    \begin{enumerate}
        \item Csak alapos üzemeltetői rutinnal rendelkező Rendszergazda töltheti be a pozíciót.

        \item Feladata a többi Rendszergazda közötti kommunikációt elősegíteni és fenntartani.

        \item Feladata a Kör által üzemeltetett rendszerek és szolgáltatások minőségét megőrizni, azokat fejleszteni az adott rendszerért felelős Rendszergazdával együttműködve.

        \item Feladata a rendszerekkel, projektekkel és szolgáltatásokkal kapcsolatban felmerülő stratégiai kérdéseket egyeztetni a Rendszergazdákkal.

        \item Rendszergazdai pozíciót szükség esetén jogosult átvenni. Ilyen esetben elszámolási kötelezetség is terheli a Körgyűlés felé.

        \item Jelentős szolgáltatáskiesés esetén, értesítését követően feladata felügyelni a helyreállítási folyamatot.

        \item A Körvezető, a Körgyűlés, a Reszortvezető vagy a Reszortgyűlés kérésére köteles beszámolni a szerverek és szolgáltatások állapotáról.

        \item A Kör a VRG-t évenként egy évre, a 6.§ által meghatározott módon választja.

        \item Kötelessége vezetőváltás esetén az új Vezető Rendszergazdát megismertetni kötelezettségeivel.
    \end{enumerate} % end of vrg

  \item Rendszergazda
    \begin{enumerate}
      \item Csak az adott feladatnak megfelelő tudással rendelkező körtag lehet.

      \item Feladata a rábízott szolgáltatás(ok), és a hozzá(juk) kapcsolódó szerver(ek) és eszközök adminisztrálása, karbantartása és lehetőségeihez mért javítása.

      \item A Vezető Rendszergazda jogosult a posztra jelentkező tudását ellenőrizni.

      \item Felel a szerver(ek)en tárolt bizalmas adatok biztonságos kezeléséért,       figyelembe véve az ide vonatkozó adatvédelmi törvényeket és egyéb      szabályozásokat.

      \item A rábízott szolgáltatás(oka)t lehetőleg napi, de legalább heti szinten      ellenőrzi.

      \item A rábízott szolgáltatások(ok)ról minél részletesebb dokumentációt készít, azt a kör tagjai számára elérhetővé teszi és azt naprakészen tartja.

      \item A rábízott szolgáltatás(ok) megfelelő időközönkénti biztonsági mentéseit megszervezi és felügyeli.

      \item Kötelessége a kör gyűlésein való részvétel.

      \item Tartja a többi Rendszergazdával, és a felügyeletére bízott eszközökhöz      kapcsolódó személyekkel a kapcsolatot.

      \item Törekszik utód kinevelésére.

      \item Jogosult belépési engedélyt kérni a szerverteremhez és a feladatainak ellátásához szükséges helyiségekhez és szolgáltatásokhoz.

      \item A rábízott szolgáltatás(ok)hoz való vészhelyzeti hozzáférést biztosító jelszavakat és más adatokat köteles a Körgyűlés által meghatározott biztonságos módon tárolni, aktuálisan tartani és a hozzáférést a Kör számára biztosítani.

      \item A rábízott szerver(eke)t köteles a Körgyűlés által meghatározott közös rendszerekkel integrálni, a rendszerek együttműködését felügyelni.

      \item A Rendszergazdákat a Körgyűlés választja, egy évre, az általuk benyújtott pályázatok alapján.
    \end{enumerate} % end of vrg

  \item Segédrendszergazda
    \begin{enumerate}
        \item Az adott feladatban minimálisan jártas személy lehet azzal a céllal, hogy
        ismereteit bővítse az adott téren.

        \item A poszt önkéntes jelentkezés vagy pályáztatás alapján kerülhet kiosztásra.

        \item A Segédrendszergazdák kiválasztása a kérdéses szerver Rendszergazdájának
        hatásköre.

        \item Feladata egy Rendszergazda segítése, apróbb feladatok elvégzése.

        \item Kötelessége a Rendszergazda iránymutatása alapján dolgozni.
    \end{enumerate} % end of Segédrendszergazda

  \item Újonc
    \begin{enumerate}
        \item A reszort által újoncnak tekintett, a körbe jelentkezett, felvett tagsággal még nem bíró személy.
    \end{enumerate}

  \item Körtag
    \begin{enumerate}
        \item Az lehet, aki felvételt nyert a Reszortba és a kör tevékenységében aktívan részt vállal.
    \end{enumerate}

\end{enumerate} % end of posztok

% +-+-+-+-+-+ posztok vége, de még felépítés

\item A Körtagság megszűnhet

  \begin{enumerate}
    \item kilépéssel,
    \item egyesületi tagság megszűnésével,
    \item KSZK tagság megszűnésével,
    \item öregtaggá válással,
    \item ha a Felügyelő Bizottság, az Elnökség, a Reszortvezetők Tanácsa vagy a Kör a 6.§ által meghatározott módon úgy dönt.
  \end{enumerate}

\end{enumerate} % felépítés vége

% Section border ==========================================================

\section{A Kör vezetője}

\begin{enumerate}
\item A Kör vezetőjét évente a Kör egy évre választja a 6.§ által meghatározott módon.
\item Feladata az újoncok munkába való bekapcsolódásának segítése.
\item Feladata a kör közösségi egységének ápolása.
\item Minden tavaszi félévben kiírja a tisztújítást és a Körnek elérhetővé teszi a Körvezetői, Rendszergazdai és Vezető Rendszergazdai posztokra érkező pályázatokat.
\item Felelős a Kör eszközeinek állapotmegóvásáért.
\item Kötelessége a Körtagokat tájékoztatni az Egyesületi döntésekről és változásokról.
\item Kötelessége vezetőváltás esetén az új Körvezetőt megismertetni kötelezettségeivel.
\item Feladata a szorgalmi időszakban lehetőleg heti rendszerességgel összehívni a Körgyűlést.
\item Feladata a Körgyűlés pontos időpontjáról, helyéről és a napirendi pontjairól értesíteni a kör minden tagját és újoncát, a Körgyűlés időpontja előtt legalább 24 órával.
\item Feladata a Körgyűlés levezetése és a napirendi pontok előterjesztése.
\end{enumerate}

% Section border ==========================================================

\section{A Kör gazdálkodása}

\begin{enumerate}
\item A Körgazdaságist a Körvezető kéri fel.
\item A Körgazdaságis a Reszortgazdaságissal együttműködve végzi munkáját.
\item A Körgazdaságist a Reszortgazdaságis kötelezheti a költségvetésről való beszámolásra.
\item A Körgazdaságis feladatai:
\begin{enumerate}
  \item a Kör gazdasági pályázatának megtervezése és elkészítése a Kör érdekeinek megfelelően,
  \item a Kör évközi költéseinek felügyelete,
  \item a Kör éves gazdasági beszámolóinak megírása.
\end{enumerate}
\end{enumerate}

% Section border ==========================================================

\section{Döntéshozatal a Körben}

\begin{enumerate}
\item A Sysadmin döntéshozó fóruma a Körgyűlés, amelyen minden felvett Körtag alanyi jogon vehet részt.
\item A Körgyűlés levezetése, a napirendi pontok előterjesztése a Körvezető feladata.
\item A Körgyűlések nyitottak, bárki részt vehet rajtuk, azonban bármikor zártkörűvé tehetőek, ha a Körgyűlés úgy határoz. Zártkörű Körgyűlésen csak a kör felvett tagjai, illetve a Körgyűlés által felkért személyek vehetnek részt.
\item A Körgyűlés határozatképes, ha a Kör tagjainak legalább fele jelen van.
\item Ha a Körgyűlés nem határozatképes, ismételt Körgyűlést kell összehívni tetszőleges határidővel azonos napirenddel, ekkor a résztvevők számától függetlenül határozatképes.
\item A felvett tagok egy szavazattal rendelkeznek, ami nem ruházható át.
\item A Kör aktív tagjainak legalább kétharmados többsége szükséges:
  \begin{enumerate}
    \item SZMSZ módosításhoz,
    \item Tag kizárásához,
    \item Poszt lemondatásához,
    \item Kör megszüntetéséhez.
  \end{enumerate}
\item Amennyiben az SZMSZ máshogy nem rendelkezik, a Körgyűlés a határozatait egyszerű többséggel hozza.
\end{enumerate}

% Section border ==========================================================

\section{A kör felvételi irányelvei}

\begin{enumerate}
\item A Körbe való felvétel és a tagság megszűnésének módját a KSZK reszort SZMSZ-e és annak mellékletei szabályozzák.
\end{enumerate}

% Section border ==========================================================

\section{Záró rendelkezések}

\begin{enumerate}
\item Jelen szabályzat által nem érintett kérdésekben a Reszort SZMSZ-e, az Egyesület Alapszabálya, valamint más vonatkozó jogszabályok és rendeletek az irányadók.
\item Jelen szabályzatot a Reszortvezetők Tanácsa fogadja el, a Kör vezetőjének javaslatára.
\item A szabályzat a Reszortvezetők Tanácsának megszavazása által lép hatályba. Ezzel egyidőben hatályát veszíti a Kör korábban elfogadott szabályzata.
\item Jelen Szervezeti és Működési Szabályzatot a Kör tagjai akaratukkal megegyezőként fogadják el.
\end{enumerate}

%=================================================================
%                           End Document
%=================================================================
\end{document}
